\begin{abstract}
\noindent
Deciding whether a desired change in cell state can be produced by a given class of
perturbation---and, if not, saying so before the experiment is run---is the central
question of cell-state engineering and of pharmaceutical target selection. Existing
computational methods either predict the response to a specified perturbation or rank
candidate regulators; none returns a \emph{feasibility verdict} grounded in measured
effects, and none can certify that a goal is unreachable. We introduce a convex-cone
reachability oracle that treats a genome-scale CRISPRi Perturb-seq screen as a dictionary
of measured effect vectors and asks whether a target cell-state shift lies in their
non-negative combination. Because loss-of-function knockdowns combine additively and
non-negatively, reachability is membership in a convex cone, decided by non-negative least
squares; a reachable target yields a minimal ranked knockdown recipe, and an unreachable
target yields a Farkas/KKT separating-hyperplane certificate that names the genes which
would instead have to be \emph{activated}---a falsifiable gain-of-function hypothesis rather
than a null result. Applied to a screen in primary human CD4\textsuperscript{+} T cells,
the oracle finds the T\textsubscript{H}2$\rightarrow$T\textsubscript{H}1 polarization shift
partly reachable by knockdown (held-out cosine $0.448$ against a shuffled-target null of
$-0.005$, $z\approx24$), with a signed decomposition attributing 39\% of the shift to
loss-of-function, 25\% to required activation, and 35\% to an irreducible remainder;
canonical regulators appear as expected positive controls (GATA3 reached, TBX21
correctly anti-aligned). Across a 12-cell atlas of four transitions under three activation
states, knockdown is never the majority modality (mean loss-of-function fraction $0.34$).
The same operator transfers unchanged to a K562 CRISPRa combinatorial screen (held-out
CEBPA cosine $0.878$), whose measured double perturbations show that the additivity
assumption is a bounded approximation (median cosine $0.71$ to the sum of singles) rather
than an identity. Against a 92-method survey (2011--2026), measured-effect grounding and a
feasibility verdict with a certificate are held by disjoint sets of prior methods; their
intersection is empty. The verdict is stable under stress: it strengthens rather than weakens
when the metric is reweighted by target magnitude to guard against signal dilution, the
non-negativity constraint costs little in fit yet supplies the certificate, and a
cross-cell-type probe shows that the \emph{direction} of a shift transfers across cell types
while the specific minimal \emph{recipe} does not. The oracle is a hypothesis-prioritization and experiment-triage
instrument---it decides where a screen should spend its arms and where it should stop---not
a target-validation engine, and every nomination is a falsifiable wet-lab hypothesis.
\end{abstract}
